\chapter{The reader base class}
\label{ch:reader}

\chapterbrief
  {Send a frame to a card and get its answer back.}
  {Abstract classes, and the Template Method pattern.}
  {\code{isCardPresent()} detecting a real DESFire.}

\section{The problem}

You can read registers. Now you need to \emph{talk to a card}, which means a
fixed sequence: tell the chip how many bits are in the last byte, flush the
FIFO, clear old interrupt flags, load the outgoing bytes, start the transceive
command, wait for a reply or a timeout, check the error register, and read the
reply out.

Here is the thing: \textbf{those steps are RC522-specific, but the protocol
above them is not.} ``Send REQA, expect a two-byte ATQA'' is ISO-14443 ---
identical on a PN532, a PN5180, or an ACR122. Write the protocol logic
\emph{inside} an RC522 class and you must rewrite all of it to support a second
chip.

\section{The idea: leave a hole in the middle of the algorithm}
\index{Template Method}

Write the protocol once, in a base class, and leave the chip-specific step
\textbf{blank} --- a method with a name and a signature but no body. Subclasses
fill the blank. The base class calls the blank as if it were already written.

That is the \textbf{Template Method} pattern: the parent owns the \emph{shape}
of the algorithm, the children supply the \emph{steps}.

\begin{analogy}{The recipe with one variable step}
A stir-fry recipe: heat the oil, \textbf{cook the protein}, add vegetables, add
sauce, serve. The sequence is fixed --- you never add sauce before the oil is hot.

But ``cook the protein'' is a hole. Chicken takes six minutes, prawns two, tofu
wants pressing first. The cookbook writes the sequence; the specific dish fills
the hole.

Critically, the recipe \emph{calls} the hole without knowing what goes in it.
\code{isCardPresent()} calls \code{transceive()}, and at runtime that call lands
in \code{Rc522Reader}, or \code{Pn532Reader}, or a fake. The base class never
learns which.
\end{analogy}

\begin{center}
\begin{tikzpicture}[font=\footnotesize\sffamily,
  tstep/.style={draw=brewergrey, rounded corners=2pt, minimum width=44mm,
               minimum height=6.5mm, inner sep=3pt},
  hole/.style={draw=breweramber, dashed, fill=breweramber!8, rounded corners=2pt,
               minimum width=44mm, minimum height=6.5mm, inner sep=3pt}]
  \node[tstep] (s1) {1. send REQA (7-bit frame)};
  \node[hole, below=2.2mm of s1] (h1) {\itshape transceive() -- the blank};
  \node[tstep, below=2.2mm of h1] (s2) {2. is the reply a valid ATQA?};
  \node[tstep, below=2.2mm of s2] (s3) {3. if silent, retry with WUPA};
  \node[draw=brewernavy, rounded corners=2pt, fit=(s1)(s3), inner sep=4mm] (box) {};
  \node[font=\scriptsize\sffamily\bfseries, color=brewernavy, above=0.5mm of box.north west,
        anchor=west] {\code{Iso14443Reader.isCardPresent()} -- final};
  \node[draw=breweramber, rounded corners=2pt, right=16mm of h1,
        minimum width=32mm, minimum height=9mm, inner sep=3pt,
        align=center, font=\scriptsize\sffamily] (impl)
        {\code{Rc522Reader}\\ fills the blank};
  \draw[-{Stealth[length=4pt]}, breweramber, dashed] (impl) -- (h1);
\end{tikzpicture}
\end{center}

\subsection{Why an abstract class and not an interface}
\index{abstract class}

An interface holds method signatures but no state. Here you need \emph{both}:
real shared code (the protocol) \emph{and} a blank (the chip step). That mix is
exactly what an abstract class is for.

\begin{center}\small
\begin{tabular}{@{}lcc@{}}
\toprule
 & \thead{interface} & \thead{abstract class} \\
\midrule
Method signatures    & yes & yes \\
Real method bodies   & mostly no & yes \\
Fields / state       & no  & yes \\
Can be instantiated  & no  & no  \\
A class can have many & yes & no (one parent) \\
\bottomrule
\end{tabular}
\end{center}

\begin{concept}{Which one to reach for}
\textbf{Interface} = ``can do this'' --- a capability, like \code{SpiLink}.

\textbf{Abstract class} = ``is a kind of this'' --- a family sharing machinery,
like \code{Iso14443Reader}.
\end{concept}

\section{Build it}

\subsection{The base}

\begin{jcode}
package com.midnightbrewer.hardware;

/**
 * An abstract base for any ISO-14443 reader.
 *
 * The protocol logic is a standard: identical for an RC522, a PN532, anything.
 * That logic lives HERE, written once. What differs per chip is only how you
 * physically push a frame out and read the reply -- left BLANK for each chip's
 * subclass to fill.
 */
public abstract class Iso14443Reader {

    // ---- ISO 14443-3 constants (chip-independent) ----
    private static final byte REQA = 0x26; // "any IDLE card, announce yourself"
    private static final byte WUPA = 0x52; // "any IDLE *or* HALTed card, wake up"

    /** A valid ATQA answer is exactly two bytes. */
    private static final int ATQA_LENGTH = 2;

    /** REQA/WUPA are 7-bit short frames, not whole bytes. */
    private static final int REQA_BITS = 7;

    // =================================================================
    // THE BLANK. 'abstract' = declared, no body. A subclass MUST supply
    // one or it will not compile.
    //
    //   sendData       : bytes to transmit
    //   bitsInLastByte : valid bits in the final byte (7 for REQA, 0 = all 8)
    //   returns        : bytes received, or an EMPTY array if nothing answered.
    //
    // Note "empty array" rather than an exception: a timeout is NORMAL when no
    // card is on the antenna. Exceptions are for the exceptional.
    // =================================================================
    protected abstract byte[] transceive(byte[] sendData, int bitsInLastByte)
            throws SpiException;

    /**
     * The second blank: compute the 2-byte CRC_A over data.
     *
     * SELECT and RATS need a CRC appended; anticollision does not. The RC522
     * has a hardware CRC unit, another chip might not -- so, like transceive,
     * it is left for the subclass. An abstract class may carry as many blanks
     * as the design needs.
     *
     * @return two CRC bytes, low first, ready to append.
     */
    protected abstract byte[] calculateCrc(byte[] data) throws SpiException;

    // =================================================================
    // THE TEMPLATE METHOD -- the fixed algorithm, shared by all chips.
    //
    // 'final' on purpose: a subclass may fill the blanks, but may NOT rewrite
    // the sequence. The protocol is the protocol.
    // =================================================================
    public final boolean isCardPresent() throws SpiException {
        // TODO 1: try REQA; if that finds a card, return true.
        // TODO 2: otherwise try WUPA and return that result.
        //
        //   Why both? REQA only wakes cards in IDLE. A card you already talked
        //   to sits in HALT and ignores REQA forever -- WUPA wakes those too.
        //   Without this, the reader works ONCE and then goes blind until the
        //   card is physically removed and replaced.
        return false;
    }

    /** Sends one REQA/WUPA and reports whether a well-formed ATQA came back. */
    private boolean request(byte mode) throws SpiException {
        // TODO: transceive a single-byte frame { mode } with REQA_BITS bits,
        //       and return whether the reply length == ATQA_LENGTH.
        return false;
    }
}
\end{jcode}

\begin{concept}{\texttt{public}, \texttt{protected}, \texttt{private}}
\code{protected} means ``visible to this class and its subclasses, but not the
outside world''. \code{transceive} is an implementation detail of the reader
family --- the kiosk should never call it.

\textbf{\code{public} for the vocabulary, \code{protected} for the family,
\code{private} for the class alone.} Choosing these deliberately is most of what
designing a class \emph{is}.

And \code{final} on \code{isCardPresent()} encodes a design decision in the type
system: \emph{chips vary, the standard does not.} Deciding what subclasses may
and may not change is the real work of writing a base class.
\end{concept}

\subsection{The RC522 subclass}

\begin{jcode}
package com.midnightbrewer.hardware;

/**
 * The RC522's answer to the two blanks in Iso14443Reader.
 *
 * 'extends' = inheritance: this IS-A Iso14443Reader and gets all its protocol
 * code for free. It supplies only the two chip-specific steps.
 */
public class Rc522Reader extends Iso14443Reader {

    // ---- registers this needs ----
    private static final int COMMAND_REG     = 0x01;
    private static final int COMM_IEN_REG    = 0x02; // which IRQs are enabled
    private static final int COMM_IRQ_REG    = 0x04; // which IRQs have fired
    private static final int ERROR_REG       = 0x06;
    private static final int FIFO_DATA_REG   = 0x09; // the FIFO, in and out
    private static final int FIFO_LEVEL_REG  = 0x0A; // bytes currently in it
    private static final int BIT_FRAMING_REG = 0x0D; // TxLastBits + StartSend

    private static final int PCD_IDLE       = 0x00; // "cancel / do nothing"
    private static final int PCD_TRANSCEIVE = 0x0C; // "send the FIFO, receive"

    // ---- CRC coprocessor ----
    private static final int DIV_IRQ_REG      = 0x05; // CRCIrq is bit 2
    private static final int CRC_RESULT_REG_L = 0x22;
    private static final int CRC_RESULT_REG_H = 0x21;
    private static final int PCD_CALCCRC      = 0x03;

    // COMPOSITION again: the reader HAS-A driver. Note it does NOT extend
    // Rc522Driver -- a reader is not a kind of driver, it USES one.
    private final Rc522Driver driver;

    public Rc522Reader(Rc522Driver driver) {
        this.driver = driver;
    }

    @Override
    protected byte[] transceive(byte[] sendData, int bitsInLastByte)
            throws SpiException {

        // ---- 1. how many bits are valid in the last byte ----
        //   C: Write(BitFramingReg, 0x07);   // 7 for REQA
        // TODO: writeRegister(BIT_FRAMING_REG, bitsInLastByte & 0x07);

        // ---- 2. arm the chip and empty the FIFO ----
        //   C: Write(CommIEnReg, 0x77|0x80);   // enable IRQs
        //      ClearBitMask(CommIrqReg, 0x80); // clear old flags
        //      SetBitMask(FIFOLevelReg, 0x80); // flush FIFO
        //      Write(CommandReg, PCD_IDLE);    // cancel running command
        // TODO: those four calls.

        // ---- 3. load the outgoing bytes into the FIFO ----
        // TODO: a loop writing each byte of sendData to FIFO_DATA_REG
        //       (as an int: sendData[i] & 0xFF).

        // ---- 4. run it, and press "start sending" ----
        //   C: Write(CommandReg, PCD_TRANSCEIVE);
        //      SetBitMask(BitFramingReg, 0x80);   // StartSend = bit 7
        // TODO: those two.

        // ---- 5. wait for the answer, or a timeout ----
        //   The C spins a large counter. That does NOT port -- see the note
        //   below. Use a WALL-CLOCK deadline:
        //
        //   long deadline = System.nanoTime() + 50_000_000L;   // ~50 ms
        //   int irq;
        //   do {
        //       irq = driver.readRegister(COMM_IRQ_REG);
        //   } while (System.nanoTime() < deadline
        //            && (irq & 0x01) == 0      // TimerIRq -> nobody answered
        //            && (irq & 0x30) == 0);    // RxIRq|IdleIRq -> done
        // TODO: the loop. Keep 'irq' -- step 6 needs it.

        // ---- 6. stop sending, then check for trouble ----
        //   C: ClearBitMask(BitFramingReg, 0x80);
        //      if (Read(ErrorReg) & 0x1B) ...   // overflow, parity, CRC, protocol
        //   Also: if TimerIRq fired, nobody answered.
        //   Either way return an EMPTY array -- "no card" is not an error.
        // TODO.

        // ---- 7. read the reply out of the FIFO ----
        //   C: n = Read(FIFOLevelReg);
        //      for (i=0;i<n;i++) backData[i] = Read(FIFODataReg);
        // TODO: read the level into n, allocate byte[n], loop reading
        //       FIFO_DATA_REG into each slot, return it.
        //
        //   Note: you read the SAME register n times. It is a FIFO -- each
        //   read pops the next byte. This looks wrong and is right.
        return new byte[0];
    }

    @Override
    protected byte[] calculateCrc(byte[] data) throws SpiException {
        //   C: Write(CommandReg, PCD_IDLE);
        //      ClearBitMask(DivIrqReg, 0x04);   // clear CRCIrq
        //      SetBitMask(FIFOLevelReg, 0x80);  // flush FIFO
        //      for (i..) Write(FIFODataReg, data[i]);
        //      Write(CommandReg, PCD_CALCCRC);
        //      wait for DivIrqReg bit 2
        //      out[0] = Read(CRCResultRegL);  out[1] = Read(CRCResultRegH);
        //
        // TODO all of it. Same wall-clock trick (~20 ms is plenty).
        // Return { low, high }: ISO-14443 puts the CRC on the wire low byte
        // first, and returning it wire-ready means callers just append it.
        return new byte[2];
    }
}
\end{jcode}

\section{The two ports that do not survive translation}

\begin{trap}{Iteration counts become deadlines}
The single most important lesson in porting embedded C to Linux.

\code{for (i = 0; i < 2000; i++)} on bare metal is microseconds. Here, each
iteration is a \code{readRegister}, which is a syscall through the kernel --- on
the order of 50\,\textmu s. The same loop count runs for \emph{minutes}.

Always time out on \code{System.nanoTime()}, never on a loop counter. If your
program appears to hang, this is the first thing to check.
\end{trap}

Second: \textbf{chip-select is no longer yours}. The C toggles CS around each
byte. Your \code{SpiLink.transfer()} already performs one complete, atomic
transaction. Try to manage CS yourself and you will fight the kernel driver and
lose.

\section{Verify on hardware}

Before polling, the chip needs a full init sequence --- the datasheet defaults
are not usable:

\begin{jcode}
private static void initChip(Rc522Driver driver) throws SpiException {
    driver.writeRegister(0x2A, 0x8D); // TModeReg     } response timer,
    driver.writeRegister(0x2B, 0x3E); // TPrescalerReg} ~300 ms at 6.78 MHz
    driver.writeRegister(0x2C, 0x02); // TReloadRegH  } (0x0258 = 600 ticks)
    driver.writeRegister(0x2D, 0x58); // TReloadRegL  }
    driver.writeRegister(0x15, 0x40); // TxASKReg  -> force 100% ASK modulation
    driver.writeRegister(0x11, 0x3D); // ModeReg   -> CRC preset 0x6363
    driver.writeRegister(0x26, 0x08); // RFCfgReg  -> receiver gain
    driver.antennaOn();
}
\end{jcode}

Two of those matter enormously:

\begin{itemize}
  \item \reg{TxASKReg} $=$ \hex{0x40} forces 100\,\% ASK modulation, which
        Type~A cards require. Without it the card never sees a valid signal, and
        you get a silence that looks exactly like a wiring fault.
  \item \reg{ModeReg} $=$ \hex{0x3D} sets the CRC preset to \hex{0x6363}, the
        ISO-14443-A value. Wrong preset means every CRC you compute is wrong,
        so every SELECT is rejected --- in Chapter~\ref{ch:activation}, not here,
        which makes it maddening to trace.
\end{itemize}

\begin{jcode}
Rc522Reader reader = new Rc522Reader(driver);
System.out.println("Polling 15x (put a card on the antenna)...");
int hits = 0;
for (int i = 0; i < 15; i++) {
    if (reader.isCardPresent()) hits++;
    Thread.sleep(40);
}
System.out.printf("card detected in %d / 15 polls%n", hits);
\end{jcode}

\begin{verifybox}
\begin{trace}
card detected in 15 / 15 polls
\end{trace}
If you get \code{0/15} with a card sitting on the antenna, check the init
sequence first --- especially those two registers --- then confirm the card is
actually 13.56\,MHz.
\end{verifybox}

\section{Traps}

\begin{traptable}
Runs for minutes, no output & Ported the C's iteration counter. Use a \code{nanoTime} deadline. \\
\code{0/15} polls, chip version fine & Missing \reg{TxASKReg}$=$\hex{0x40} or \reg{ModeReg}$=$\hex{0x3D}. \\
Detects once, then never again & Only sending REQA. Halted cards need WUPA. \\
Detection works about half the time & Wire length, or card not flat on the antenna. \\
\code{NegativeArraySizeException} & FIFO level read before the reply arrived --- check the wait loop. \\
\end{traptable}

\section{Commit}

\begin{shell}
git add -A
git commit -m "M3: Iso14443Reader template method + Rc522Reader transceive/CRC"
\end{shell}
